\section{Referencias Bibliográficas}
\label{sec:referencias}

% Formato APA con sangria francesa (hanging indent), tal como en el
% informe original. No se usa biblatex/biber: las citas dentro del
% texto van escritas literalmente como texto parentetico, igual que en
% el documento fuente (p. ej. "(Li, Zhang y Chen, 2018)").
\newcommand{\refapa}[1]{\par\noindent\hangindent=1.25cm\hangafter=1 #1\par\vspace{0.35cm}}

% \url no puede escribirse directamente dentro del argumento de \refapa cuando
% contiene caracteres especiales de TeX (%, _, #, ~, ^, &): ese argumento se
% tokeniza ANTES de que \url cambie los codigos de categoria, asi que "_" abre
% subindice y "%" abre un comentario que se traga el resto de la linea (el
% error "File ended while scanning use of \refapa" viene de ahi). \urldef
% arma el enlace por separado, en la categoria de codigos correcta, y aqui
% solo se referencia el nombre ya seguro.
\urldef{\urlConstitucion}\url{https://leyes.congreso.gob.pe/Documentos/constituciones_ordenado/CONSTIT_1993/Texto_actualizado_CONS_1993.pdf}

\refapa{Congreso de la República del Perú. (1993). \emph{Constitución Política del Perú},
artículo 2. Texto actualizado: \urlConstitucion}

\refapa{Congreso de la República del Perú. (2011). \emph{Ley N.º 29733, Ley de Protección
de Datos Personales}. Diario Oficial El Peruano.
\url{https://www.gob.pe/institucion/congreso-de-la-republica/normas-legales/243470-29733}}

\refapa{Congreso de la República del Perú. (2013). \emph{Ley N.º 30120, Ley de apoyo a la
seguridad ciudadana con cámaras de videovigilancia públicas y privadas}. Diario Oficial El
Peruano. \url{https://www.leyes.congreso.gob.pe/Documentos/Leyes/30120.pdf}}

\urldef{\urlDirectiva}\url{https://www.minjus.gob.pe/wp-content/uploads/2020/01/Directiva-N°-01-2020-DGTAIPD-1.pdf}

\refapa{Dirección General de Transparencia, Acceso a la Información Pública y Protección
de Datos Personales. (2020). \emph{Directiva N.º 01-2020-JUS/DGTAIPD, para el tratamiento
de datos personales mediante sistemas de videovigilancia}. Ministerio de Justicia y Derechos
Humanos. \urlDirectiva}

\refapa{Ministerio de Justicia y Derechos Humanos. (2024). \emph{Decreto Supremo N.º
016-2024-JUS, que aprueba el Reglamento de la Ley N.º 29733, Ley de Protección de Datos
Personales}. Diario Oficial El Peruano.
\url{https://www.gob.pe/institucion/smv/normas-legales/6426760-016-2024-jus}}

\refapa{Ministerio del Interior. (2020). \emph{Decreto Supremo N.º 007-2020-IN, que
aprueba el Reglamento del Decreto Legislativo N.º 1218 y de la Ley N.º 30120}. Diario Oficial
El Peruano.
\url{https://www.gob.pe/institucion/mininter/normas-legales/1081739-007-2020-in}}

\refapa{Poder Ejecutivo. (2015). \emph{Decreto Legislativo N.º 1218, que regula el uso de las
cámaras de videovigilancia}. Diario Oficial El Peruano.
\url{https://busquedas.elperuano.pe/normaslegales/decreto-legislativo-que-regula-el-uso-de-las-camaras-de-vide-decreto-legislativo-n-1218-1291565-8/}}

\bigskip
\noindent\textbf{Trabajos técnicos}
\medskip

\refapa{Aharon, N., Orfaig, R., \& Bobrovsky, B.-Z. (2022). BoT-SORT: Robust associations
multi-pedestrian tracking. \emph{arXiv preprint arXiv:2206.14651}.
\url{https://arxiv.org/abs/2206.14651}}

\refapa{Dalal, N., \& Triggs, B. (2005). Histograms of oriented gradients for human detection.
En \emph{Proceedings of the IEEE Computer Society Conference on Computer Vision and Pattern
Recognition} (pp. 886-893). IEEE. \url{https://doi.org/10.1109/CVPR.2005.177}}

\refapa{Dendorfer, P., Rezatofighi, H., Milan, A., Shi, J., Cremers, D., Reid, I., Roth, S.,
Schindler, K., \& Leal-Taixé, L. (2020). MOT20: A benchmark for multi object tracking in crowded
scenes. \emph{arXiv preprint arXiv:2003.09003}. \url{https://arxiv.org/abs/2003.09003}}

\refapa{Jocher, G., \& Qiu, J. (2024). \emph{Ultralytics YOLO11} (versión 11.0.0) [Software].
Ultralytics. \url{https://github.com/ultralytics/ultralytics}}

\refapa{Li, Y., Zhang, X., \& Chen, D. (2018). CSRNet: Dilated convolutional neural
networks for understanding the highly congested scenes. En \emph{Proceedings of the IEEE
Conference on Computer Vision and Pattern Recognition} (pp. 1091-1100). IEEE.
\url{https://arxiv.org/abs/1802.10062}}

\refapa{Lin, W., Zhao, C., \& Chan, A. B. (2025). Point-to-region loss for semi-supervised
point-based crowd counting. En \emph{Proceedings of the IEEE/CVF Conference on Computer
Vision and Pattern Recognition} (pp. 29363-29373). IEEE.
\url{https://arxiv.org/abs/2505.21943}}

\refapa{Redmon, J., Divvala, S., Girshick, R., \& Farhadi, A. (2016). You only look once:
Unified, real-time object detection. En \emph{Proceedings of the IEEE Conference on Computer
Vision and Pattern Recognition} (pp. 779-788). IEEE. \url{https://arxiv.org/abs/1506.02640}}

\refapa{Song, Q., Wang, C., Jiang, Z., Wang, Y., Tai, Y., Wang, C., Li, J., Huang, F., \& Wu,
Y. (2021). Rethinking counting and localization in crowds: A purely point-based framework.
En \emph{Proceedings of the IEEE/CVF International Conference on Computer Vision}
(pp. 3365-3374). IEEE. \url{https://arxiv.org/abs/2107.12746}}

\refapa{Zhang, Y., Sun, P., Jiang, Y., Yu, D., Weng, F., Yuan, Z., Luo, P., Liu, W., \& Wang,
X. (2022). ByteTrack: Multi-object tracking by associating every detection box. En
\emph{Proceedings of the European Conference on Computer Vision} (pp. 1-21). Springer.
\url{https://arxiv.org/abs/2110.06864}}

\refapa{Zhou, K., Yang, Y., Cavallaro, A., \& Xiang, T. (2019). Omni-scale feature learning
for person re-identification. En \emph{Proceedings of the IEEE/CVF International Conference
on Computer Vision} (pp. 3702-3712). IEEE. \url{https://arxiv.org/abs/1905.00953}}
